\chapter{Analiza wyników przeprowadzonych badan}
Program kompilatora jest niezbędnym narzędziem pracy, dla każdgo inżyniera 
wytwarzającego oprogramowanie, który zwraca uwagę na zużycie zasobów maszyny wykonującej jego program.
W wielu przypadkach narzędzie to jest traktowane jako czarna
skrzynka - kod źródłowy jest wejściem, na wyjściu otrzymujemy plik wykonywalny , a to, co dzieje
się pomiędzy, pozostaje poza uwagą programisty. Przy próbie optymalizacji programu zwykle to logika i przebieg algorytmu 
jest modyfikowana jako pierwsza, pomijając wpływ, jaki na wydajność ma proces tłumaczenia kodu na instrukcje maszynowe.

Celem przeprowadzonych badań było sprawdzenie na ile możliwa jest poprawa wydajności implementacji FFT z użyciem różnych zabiegów.
Jednym z nich była weryfikacja ile daje sama zmiana flag optymalizacji kompialatora przy niezmienionym
kodzie algorytmu. Weryfikowano również czy jawne użycie instrukcji wykorzystania sprzętu fizycznego, nawet w ograniczonej skali takiej jak 
zmiana układu danych wejściowych czy ręczne użycie instrukcji wektorowych, może dać realny wpływ na optymalizacje. 

Wyniki i wnioski z przeprowadzonych badań mogą mieć aspekt praktyczny, zwłaszcza w systemach wbudowanych, które często mają ograniczone
zasoby. W ich przypadku większa wydajność, bez droższego sprzętu fizycznego, oznacza niższy koszt
i pobór energii,co może realnie poprawić działanie, stabilność i rentowność całego systemu.

Celem niniejszego rozdziału jest przedstawienie uzyskanych wyników oraz próba
ich wyjaśnienia poprzez analizę przyczyn stojących za zmianą poszczególnych
parametrów optymalizacyjnych.

\section{Powtarzalność pomiarów i próg istotności różnic}\label{sec:prog-istotnosci}
Przed omawianiem wyników należy zaznaczyć, że zgodnie z metodyką przedstawioną w~\cref{sec:metodykaBadan}, przeprowadzone badania dały bardzo
obszerny zbiór wyników, który w całości załączono jako załącznik do
niniejszej pracy. Zestaw ten zawiera 100 tabel, zbliżonych formą do
prezentowanych w rozdziale. Tabele przedstawione poniżej wybrano
spośród nich celowo, tak aby zilustrować najbardziej interesujące
mechanizmy zaobserwowane w trakcie badań.

Jakość środowiska pomiarowego dla wersji sekwencyjnej, z danymi wejściowymi o 1048576 próbkach  ilustruje
\cref{tab:iqr-gcc-1M}. Rozstęp względny większości wariantów nie
przekracza 3\%.
W załaczonych wynikach tabela dla wersji współbieżnej pokazuje bardzo zbliżony wzorzec.
Rozstęp względny między pomiarami nie przekracza tam 4\%, co potwierdza, że wysoka
powtarzalność pomiarów nie jest cechą specyficzną dla wariantów
sekwencyjnych, lecz efektem zastosowanej metodyki.

Na tej podstawie można pokusić się o stwierdzenie, że różnice w czasie
wykonania przekraczające 5\% wynikają z rzeczywistego wpływu flag
kompilacji lub decyzji implementacyjnych, a nie z szumu pomiarowego. Należy brać jednak pod uwagę, że w innych konfiguracjach macierzy
testów, na przykład przy mniejszych rozmiarach danych wejściowych rozstęp względny bywa wyższy (ale wciąż nie przekraczający 7\%), co należy uwzględnić przy
interpretacji różnic bliskich granicy.
%============================================================================
% Tabela 1.15   (etykieta: tab:iqr-gcc-1M)
% Rozstęp międzykwartylowy bezwzględny i względny dla GCC i 1048576 próbek, wersja sekwencyjna
%============================================================================

\begin{table}[H]
\centering
\caption{Rozstęp międzykwartylowy bezwzględny i względny dla GCC i 1048576 próbek, wersja sekwencyjna.}\label{tab:iqr-gcc-1M}
\scriptsize
\setlength{\tabcolsep}{3.0pt}
\begin{tabular}{lrrrrrrrrrrr}
\toprule
Algorytm & \texttt{O0} & \texttt{O1} & \texttt{O2} & \texttt{O3} & \texttt{Os} & unroll & ffast & mcpu & ff+mcpu & +alias & lto \\
\midrule
\multicolumn{12}{l}{\textit{Rozstęp międzykwartylowy IQR [ms]}} \\[2pt]
radix2\_v0 & 1,197 & 3,271 & 2,567 & 2,570 & 3,411 & 2,688 & 3,129 & 3,684 & 1,575 & 2,685 & 1,051 \\
radix2\_v1 & 15,923 & 1,872 & 0,540 & 1,279 & 1,603 & 1,593 & 1,415 & 1,197 & 1,634 & 1,720 & 1,336 \\
radix2\_v2\_soa & 5,861 & 2,353 & 3,665 & 2,685 & 8,621 & 3,282 & 3,947 & 2,987 & 3,155 & 2,241 & 2,272 \\
radix2\_v3\_neon & 4,969 & 5,419 & 4,200 & 4,928 & 5,496 & 6,389 & 3,962 & 1,507 & 2,235 & 3,438 & 6,759 \\
radix4 & 5,112 & 3,821 & 4,789 & 6,146 & 4,149 & 3,889 & 4,990 & 3,929 & 3,847 & 2,535 & 3,556 \\
radix4\_naive & 4,731 & 1,170 & 1,175 & 1,252 & 1,365 & 1,374 & 1,333 & 1,441 & 0,343 & 0,224 & 1,298 \\
splitradix & 11,586 & 6,939 & 5,438 & 10,603 & 12,877 & 4,317 & 7,192 & 7,439 & 4,638 & 3,148 & 11,708 \\
stockham & 4,809 & 0,706 & 0,432 & 0,729 & 1,280 & 0,223 & 0,826 & 0,952 & 0,240 & 0,493 & 0,172 \\
\addlinespace
\multicolumn{12}{l}{\textit{Rozstęp względny IQR/mediana [\%]}} \\[2pt]
radix2\_v0 & 0,14 & 1,54 & 1,24 & 1,24 & 1,63 & 1,30 & 1,52 & 1,78 & 0,76 & 1,29 & 0,50 \\
radix2\_v1 & 2,49 & 0,88 & 0,26 & 0,61 & 0,76 & 0,77 & 0,68 & 0,58 & 0,78 & 0,82 & 0,63 \\
radix2\_v2\_soa & 1,24 & 1,11 & 1,75 & 1,29 & 4,15 & 1,55 & 1,82 & 1,43 & 1,46 & 1,04 & 1,07 \\
radix2\_v3\_neon & 1,18 & 2,55 & 1,99 & 2,34 & 2,56 & 3,06 & 1,86 & 0,71 & 1,05 & 1,60 & 3,18 \\
radix4 & 0,79 & 1,26 & 1,71 & 2,20 & 1,51 & 1,47 & 1,62 & 1,45 & 1,28 & 0,84 & 1,25 \\
radix4\_naive & 0,75 & 0,46 & 0,46 & 0,42 & 0,54 & 0,56 & 0,53 & 0,55 & 0,14 & 0,09 & 0,50 \\
splitradix & 1,55 & 1,09 & 0,91 & 1,78 & 2,14 & 0,72 & 1,14 & 1,34 & 0,77 & 0,52 & 1,95 \\
stockham & 1,55 & 0,41 & 0,25 & 0,80 & 0,76 & 0,24 & 0,91 & 1,04 & 0,26 & 0,54 & 0,19 \\
\bottomrule
\end{tabular}
\end{table}

\section{Analiza wpływu optymalizacji programu sekwencyjnego}

\subsection{Wpływ poziomów optymalizacji kompilatora}
Pierwsze rozważania dotyczą wpływu stosowanych flag kompilatora na przygotowane algorytmu sekwencyjne (inaczej jednowątkowe).

Robocze:
Przeskok z O0 daje najwiecej zdecydowanie, pozniej te czasy sie wyplaszczaja. Loop unrooling widac po rozmiarach .text z asm

%============================================================================
% Tabela 1.2+1.5 -- tab:med-merged-16k
% Mediana czasu wykonania transformaty dla GCC i Clang 16384 próbek, wersja sekwencyjna. Wartości w~milisekundach.
%============================================================================
\begin{table}[H]
\centering
\caption{Mediana czasu wykonania transformaty dla GCC i Clang, 16384 próbek, wersja sekwencyjna. Wartości w~milisekundach. Pogrubiono krótszy z dwóch czasów w~każdej komórce.}
\label{tab:med-merged-16k}
\scriptsize
\setlength{\tabcolsep}{3.0pt}
\begin{tabular}{llrrrrrrrrrrr}
\toprule
Algorytm & Kompilator & \texttt{O0} & \texttt{O1} & \texttt{O2} & \texttt{O3} & \texttt{Os} & unroll & ffast & mcpu & ff+mcpu & +alias & lto \\
\midrule
\multirow{2}{*}{radix2\_v0}
  & GCC   & 8,6189 & 1,0518 & 0,9270 & 0,9421 & 0,9592 & 0,8966 & 0,9259 & 0,9315 & 0,9156 & 0,9143 & 0,9405 \\
  & Clang & \textbf{6,4565} & \textbf{0,9253} & \textbf{0,9261} & \textbf{0,9018} & \textbf{0,9199} & \textbf{0,9060} & \textbf{0,8185} & \textbf{0,7955} & \textbf{0,8049} & \textbf{0,8040} & \textbf{0,9320} \\
\addlinespace
\multirow{2}{*}{radix2\_v1}
  & GCC   & 6,4426 & 0,9004 & 0,7562 & 0,7563 & 0,9076 & 0,7775 & 0,8442 & 0,8871 & 0,8472 & 0,8453 & 0,7858 \\
  & Clang & \textbf{5,0313} & \textbf{0,6843} & \textbf{0,7063} & \textbf{0,7172} & \textbf{0,6778} & \textbf{0,7136} & \textbf{0,7277} & \textbf{0,7377} & \textbf{0,7551} & \textbf{0,7569} & \textbf{0,7188} \\
\addlinespace
\multirow{2}{*}{radix2\_v2\_soa}
  & GCC   & 4,2421 & 0,9567 & 0,8427 & 0,8594 & 0,9564 & 0,8606 & 0,8561 & 0,8834 & 0,8826 & \textbf{0,7111} & 0,9681 \\
  & Clang & \textbf{4,1100} & \textbf{0,8651} & \textbf{0,6758} & \textbf{0,6579} & \textbf{0,8528} & \textbf{0,6484} & \textbf{0,6248} & \textbf{0,6239} & \textbf{0,6290} & 0,6153 & \textbf{0,7038} \\
\addlinespace
\multirow{2}{*}{radix2\_v3\_neon}
  & GCC   & \textbf{3,6592} & \textbf{0,5794} & \textbf{0,5624} & \textbf{0,5643} & \textbf{0,6118} & \textbf{0,5635} & \textbf{0,5656} & \textbf{0,5721} & \textbf{0,5840} & \textbf{0,6496} & \textbf{0,5872} \\
  & Clang & 4,5329 & 0,6017 & 0,5995 & 0,6188 & 0,5986 & 0,6420 & 0,6603 & 0,6267 & 0,6265 & 0,6362 & 0,6335 \\
\addlinespace
\multirow{2}{*}{radix4}
  & GCC   & 4,6056 & 0,9789 & 0,9057 & 0,8940 & 0,9521 & 0,8850 & 0,9829 & 0,9307 & 0,9901 & 0,9858 & 0,9310 \\
  & Clang & \textbf{3,7740} & \textbf{0,9248} & \textbf{0,9086} & \textbf{0,8852} & \textbf{0,8792} & \textbf{0,8759} & \textbf{0,9005} & \textbf{0,8518} & \textbf{0,8478} & \textbf{0,8480} & \textbf{0,8920} \\
\addlinespace
\multirow{2}{*}{radix4\_naive}
  & GCC   & 6,0956 & \textbf{2,1677} & 3,0243 & 3,0293 & \textbf{2,2145} & 3,0125 & \textbf{2,0307} & 3,0414 & \textbf{1,9943} & \textbf{1,9947} & \textbf{2,1756} \\
  & Clang & \textbf{5,1931} & 2,7980 & \textbf{2,7485} & \textbf{2,7394} & 2,7238 & \textbf{2,6721} & 2,1736 & \textbf{2,7164} & 2,1778 & 2,1781 & 2,6662 \\
\addlinespace
\multirow{2}{*}{splitradix}
  & GCC   & \textbf{3,1690} & \textbf{1,1991} & 1,0808 & 1,0852 & 1,1302 & 1,1023 & 1,1422 & 1,0561 & 1,1259 & 1,1261 & 1,1097 \\
  & Clang & 3,2003 & 1,0451 & \textbf{1,0273} & \textbf{1,0389} & \textbf{1,0633} & \textbf{1,0200} & \textbf{0,9967} & \textbf{0,9883} & \textbf{1,0310} & \textbf{1,0090} & \textbf{1,0474} \\
\addlinespace
\multirow{2}{*}{stockham}
  & GCC   & 3,5101 & 0,6360 & 0,6060 & 0,6218 & \textbf{0,5544} & 0,6218 & 0,6022 & 0,6189 & 0,6149 & 0,6203 & 0,6245 \\
  & Clang & \textbf{3,5632} & \textbf{0,4995} & \textbf{0,4456} & \textbf{0,4475} & 0,4702 & \textbf{0,4395} & \textbf{0,4433} & \textbf{0,4742} & \textbf{0,4478} & \textbf{0,4548} & \textbf{0,4296} \\
\addlinespace
\multirow{2}{*}{fftw}
  & GCC   & 0,3596 & 0,3764 & 0,3513 & 0,3704 & 0,3649 & \textbf{0,3514} & 0,3701 & 0,3730 & \textbf{0,3498} & 0,3679 & 0,3631 \\
  & Clang & \textbf{0,3516} & \textbf{0,3535} & \textbf{0,3470} & \textbf{0,3625} & \textbf{0,3502} & 0,3570 & \textbf{0,3671} & \textbf{0,3659} & 0,3630 & \textbf{0,3622} & \textbf{0,3589} \\
\bottomrule
\end{tabular}
\end{table}



\subsection{Wpływ celowych modyfikacji kodu programu, w tym wybór algorytmu}
Robocze 
Widac, ze neon jest duzo szybszy od v1 i v2, pewnie gdyby wykorzystac tez te instrukcje neon dla pozostalych to byloby jeszcze lepiej.
Soa od poziomu 02, poza 0s jest tez szybsza troche, ale tylko dla Clanga tez w innych wariantach to wiekszy danych widac. Poziom O1 działa najgorzej dla soa, trzeba zobaczyc w ASM dlaczego.
Wyliczanie kazdorazowe w wersjach v0 i naive zabija predkosc wywolania, dla rozsadnych rozmiarow warto przekazywac wspolczynniki jako parametry wejsciowe, a nie liczyc je w locie.
%============================================================================
% Tabela 1.8   (etykieta: tab:speedup-gcc-1M)
% Przyspieszenie względem -O0 dla GCC i 16384 próbek, wersja sekwencyjna. Wartości bezwymiarowe, poziom \texttt{-O0} przyjęto za~1,00.
%============================================================================
\begin{table}[H]
\centering
\caption{Przyspieszenie względem -O0 dla GCC i 16384 próbek, wersja sekwencyjna. Wartości bezwymiarowe, poziom \texttt{-O0} przyjęto za~1,00.}
\label{tab:speedup-gcc-16k}
\scriptsize
\setlength{\tabcolsep}{3.0pt}
\begin{tabular}{lrrrrrrrrrrr}
\toprule
Algorytm & \texttt{O0} & \texttt{O1} & \texttt{O2} & \texttt{O3} & \texttt{Os} & unroll & ffast & mcpu & ff+mcpu & +alias & lto \\
\midrule
radix2\_v0 & 1,00 & 8,19 & 9,30 & 9,15 & 8,99 & 9,61 & 9,31 & 9,25 & 9,41 & 9,43 & 9,16 \\
radix2\_v1 & 1,00 & 7,15 & 8,52 & 8,52 & 7,10 & 8,29 & 7,63 & 7,26 & 7,60 & 7,62 & 8,20 \\
radix2\_v2\_soa & 1,00 & 4,43 & 5,03 & 4,94 & 4,44 & 4,93 & 4,96 & 4,80 & 4,81 & 5,97 & 4,38 \\
radix2\_v3\_neon & 1,00 & 6,32 & 6,51 & 6,48 & 5,98 & 6,49 & 6,47 & 6,40 & 6,27 & 5,63 & 6,23 \\
radix4 & 1,00 & 4,70 & 5,08 & 5,15 & 4,84 & 5,20 & 4,69 & 4,95 & 4,65 & 4,67 & 4,95 \\
radix4\_naive & 1,00 & 2,81 & 2,02 & 2,01 & 2,75 & 2,02 & 3,00 & 2,00 & 3,06 & 3,06 & 2,80 \\
splitradix & 1,00 & 2,64 & 2,93 & 2,92 & 2,80 & 2,88 & 2,77 & 3,00 & 2,81 & 2,81 & 2,86 \\
stockham & 1,00 & 5,52 & 5,79 & 5,65 & 6,33 & 5,65 & 5,83 & 5,67 & 5,71 & 5,66 & 5,62 \\
\bottomrule
\end{tabular}
\end{table}
%tutaj o instrukcjach SIMD, ulozenie danych 

\subsection{Porównanie kompilatorów GCC i Clang}
Choć oba kompilatory generują mniejszy kod z podobną częstością (GCC w 45 z
88 przypadków, Clang w 42), średni rozmiar kodu wynikowego dla GCC (518~B)
jest o 40\% mniejszy niż dla Clang (724~B). 

Clang generuje szybszy kod w 78 z 99 przypadków (79\%), przy średniej
przewadze rzędu 1,15 raza; łączny średni czas wykonania jest dla Clang
o~6,7\% niższy niż dla GCC (z~pominięciem \texttt{-O0}).
%============================================================================
% Tabela 1.53 i 1.54   (etykieta: tab:med-clang-16k)
% Rozmiar sekcji kodu wynikowego algorytmu dla GCC i Clang, wersja sekwencyjna. Rozmiar sekcj.
%============================================================================
\begin{table}[H]
\centering
\caption{Rozmiar sekcji kodu wynikowego algorytmu dla GCC i Clang, wersja sekwencyjna. Rozmiar sekcji \texttt{.text} w~bajtach. Pogrubiono mniejszą z dwóch wartości w~każdej komórce.}
\label{tab:text-merged}
\scriptsize
\setlength{\tabcolsep}{3.0pt}
\begin{tabular}{llrrrrrrrrrrr}
\toprule
Algorytm & Kompilator & \texttt{O0} & \texttt{O1} & \texttt{O2} & \texttt{O3} & \texttt{Os} & unroll & ffast & mcpu & ff+mcpu & +alias & lto \\
\midrule
\multirow{2}{*}{radix2\_v0}
  & GCC   & 732  & 344 & 348  & 344  & \textbf{276} & 580  & 344  & 336  & 336  & 336  & 344 \\
  & Clang & \textbf{628}  & \textbf{292} & \textbf{292}  & \textbf{280}  & \textbf{276} & \textbf{280}  & \textbf{272}  & \textbf{280}  & \textbf{280}  & \textbf{280}  & \textbf{280} \\
\addlinespace
\multirow{2}{*}{radix2\_v1}
  & GCC   & 704  & 260 & \textbf{224}  & \textbf{228}  & 236 & \textbf{392}  & \textbf{260}  & \textbf{244}  & \textbf{244}  & \textbf{244}  & \textbf{228} \\
  & Clang & \textbf{604}  & \textbf{236} & 840  & 852  & \textbf{224} & 852  & 848  & 864  & 860  & 860  & 852 \\
\addlinespace
\multirow{2}{*}{radix2\_v2\_soa}
  & GCC   & 784  & \textbf{272} & \textbf{268}  & \textbf{272}  & \textbf{260} & \textbf{460}  & \textbf{272}  & \textbf{264}  & \textbf{264}  & 1032 & \textbf{272} \\
  & Clang & \textbf{652}  & 288 & 712  & 740  & 264 & 740  & 736  & 772  & 768  & \textbf{768}  & 740 \\
\addlinespace
\multirow{2}{*}{radix2\_v3\_neon}
  & GCC   & 1448 & \textbf{404} & \textbf{384}  & \textbf{564}  & \textbf{360} & \textbf{832}  & \textbf{564}  & \textbf{560}  & \textbf{560}  & \textbf{1068} & \textbf{564} \\
  & Clang & \textbf{1396} & 444 & 1048 & 1536 & 412 & 1536 & 1532 & 1528 & 1528 & 1528 & 1536 \\
\addlinespace
\multirow{2}{*}{radix4}
  & GCC   & 1216 & 476 & 452  & 456  & 428 & 456  & 472  & 440  & 468  & 468  & 456 \\
  & Clang & \textbf{1068} & \textbf{416} & \textbf{416}  & \textbf{384}  & \textbf{416} & \textbf{384}  & \textbf{384}  & \textbf{396}  & \textbf{396}  & \textbf{396}  & \textbf{384} \\
\addlinespace
\multirow{2}{*}{radix4\_naive}
  & GCC   & \textbf{1188} & \textbf{536} & \textbf{512}  & \textbf{512}  & \textbf{500} & \textbf{512}  & \textbf{572}  & \textbf{524}  & \textbf{572}  & \textbf{572}  & \textbf{512} \\
  & Clang & 1080 & 560 & 560  & 548  & 548 & 548  & 2208 & 552  & 2232 & 2232 & 548 \\
\addlinespace
\multirow{2}{*}{splitradix}
  & GCC   & 1444 & 580 & 528  & 528  & 532 & 528  & 540  & 532  & 544  & 544  & 528 \\
  & Clang & \textbf{1392} & \textbf{464} & \textbf{464}  & \textbf{464}  & \textbf{448} & \textbf{464}  & \textbf{464}  & \textbf{476}  & \textbf{476}  & \textbf{476}  & \textbf{464} \\
\addlinespace
\multirow{2}{*}{stockham}
  & GCC   & 764  & 320 & \textbf{320}  & \textbf{684}  & \textbf{208} & 1544 & \textbf{620}  & \textbf{668}  & \textbf{668}  & \textbf{668}  & \textbf{684} \\
  & Clang & \textbf{716}  & \textbf{268} & 1080 & 852  & 216 & \textbf{852}  & 840  & 1092 & 868  & 868  & 852 \\
\bottomrule
\end{tabular}
\end{table}



\subsection{Wpływ kompilatora na błędy numeryczny wybranych algorytmów}\label{sec:WplywBledyNumeryczne}
%absolutnie odnies sie do tabeli 5,1 (albo 1.15)
Tutaj implementacja v0 ma blad numeryczny w wyliczaniu motylka, bylo pokazane w rozdziale wczesniej, trzeba zrobic cref tutaj. Warty do zauwazenia jest wniosek ze 
kompilator to narzedzie, a nie poprawiacz kodu - jak napiszemy cos po prostu zle to on nie poprawi magicznie dokladnosci, tylko ten blad bedzie. Ten blad numeryczny
jest tym bardziej widoczny im wieksze sa dane wejsciowe.
%============================================================================
% Tabela 1.44   (etykieta: tab:dokl-gcc-1M)
% Błąd średniokwadratowy i stosunek sygnału do szumu dla GCC i 1048576 próbek, wersja sekwencyjna.
%============================================================================
\begin{table}[H]
\centering
\caption{Błąd średniokwadratowy i stosunek sygnału do szumu dla GCC i 16384 próbek, wersja sekwencyjna.}
\label{tab:dokl-gcc-16k}
\scriptsize
\setlength{\tabcolsep}{3.0pt}
\begin{tabular}{lrrrrrrrrrrr}
\toprule
Algorytm & \texttt{O0} & \texttt{O1} & \texttt{O2} & \texttt{O3} & \texttt{Os} & unroll & ffast & mcpu & ff+mcpu & +alias & lto \\
\midrule
\multicolumn{12}{l}{\textit{Błąd średniokwadratowy RMSE}} \\[2pt]
radix2\_v0 $(\cdot 10^{-3})$ & 1,714 & 1,714 & 2,595 & 2,595 & 2,595 & 2,595 & 2,594 & 2,595 & 2,594 & 2,594 & 2,595 \\
radix2\_v1 $(\cdot 10^{-5})$ & 1,377 & 1,377 & 1,429 & 1,429 & 1,429 & 1,429 & 1,444 & 1,429 & 1,444 & 1,444 & 1,429 \\
radix2\_v2\_soa $(\cdot 10^{-5})$ & 1,377 & 1,377 & 1,429 & 1,429 & 1,429 & 1,429 & 1,444 & 1,429 & 1,444 & 1,444 & 1,429 \\
radix2\_v3\_neon $(\cdot 10^{-5})$ & 1,445 & 1,445 & 1,467 & 1,467 & 1,467 & 1,467 & 1,542 & 1,467 & 1,542 & 1,542 & 1,467 \\
radix4 $(\cdot 10^{-5})$ & 1,365 & 1,365 & 1,174 & 1,174 & 1,174 & 1,174 & 1,236 & 1,174 & 1,236 & 1,236 & 1,174 \\
radix4\_naive $(\cdot 10^{-5})$ & 1,760 & 1,760 & 1,610 & 1,610 & 1,610 & 1,610 & 1,571 & 1,610 & 1,571 & 1,571 & 1,610 \\
splitradix $(\cdot 10^{-5})$ & 1,276 & 1,276 & 1,280 & 1,280 & 1,280 & 1,280 & 1,165 & 1,280 & 1,165 & 1,165 & 1,280 \\
stockham $(\cdot 10^{-5})$ & 1,342 & 1,342 & 1,387 & 1,387 & 1,387 & 1,387 & 1,429 & 1,387 & 1,429 & 1,429 & 1,387 \\
\addlinespace
\multicolumn{12}{l}{\textit{Stosunek sygnału do szumu SNR [dB]}} \\[2pt]
radix2\_v0 & 95,12 & 95,12 & 91,52 & 91,52 & 91,52 & 91,52 & 91,53 & 91,52 & 91,53 & 91,53 & 91,52 \\
radix2\_v1 & 137,03 & 137,03 & 136,71 & 136,71 & 136,71 & 136,71 & 136,61 & 136,71 & 136,61 & 136,61 & 136,71 \\
radix2\_v2\_soa & 137,03 & 137,03 & 136,71 & 136,71 & 136,71 & 136,71 & 136,61 & 136,71 & 136,61 & 136,61 & 136,71 \\
radix2\_v3\_neon & 136,61 & 136,61 & 136,48 & 136,48 & 136,48 & 136,48 & 136,04 & 136,48 & 136,04 & 136,04 & 136,48 \\
radix4 & 137,10 & 137,10 & 138,41 & 138,41 & 138,41 & 138,41 & 137,96 & 138,41 & 137,96 & 137,96 & 138,41 \\
radix4\_naive & 134,89 & 134,89 & 135,67 & 135,67 & 135,67 & 135,67 & 135,88 & 135,67 & 135,88 & 135,88 & 135,67 \\
splitradix & 137,69 & 137,69 & 137,66 & 137,66 & 137,66 & 137,66 & 138,48 & 137,66 & 138,48 & 138,48 & 137,66 \\
stockham & 137,25 & 137,25 & 136,96 & 136,96 & 136,96 & 136,96 & 136,70 & 136,96 & 136,70 & 136,70 & 136,96 \\
\bottomrule
\end{tabular}
\end{table}

%============================================================================
% Tabela 1.47   (etykieta: tab:dokl-clang-1M)
% Błąd średniokwadratowy i stosunek sygnału do szumu dla Clang i 16384 próbek, wersja sekwencyjna.
%============================================================================
\subsection{Wpływ wybranych flag na wektoryzację}
TO DO 
!!!!!!!!!!!Ta tabela jest straszna, trzeba wymyślić jakis sensowniejszy sposób prezentacji tych wyników!!!!!
%============================================================================
% Tabela 1.51 i 1,52  (etykieta: tab:simd-skal-gcc)
% Liczba instrukcji wektorowych i skalarnych dla GCC i clang, wersja sekwencyjna. Liczba statycznych instrukcji zmiennoprzecinkowych w~kodzie wynikowym.
%============================================================================
\begin{table}[H]
\centering
\caption{Liczba instrukcji wektorowych i skalarnych dla GCC i Clang, wersja sekwencyjna. Liczba statycznych instrukcji zmiennoprzecinkowych w~kodzie wynikowym. W~sekcji wektorowej pogrubiono wyższą wartość (więcej wektoryzacji), w~sekcji skalarnej -- niższą.}
\label{tab:simd-skal-merged}
\scriptsize
\setlength{\tabcolsep}{3.0pt}
\begin{tabular}{llrrrrrrrrrrr}
\toprule
Algorytm & Kompilator & \texttt{O0} & \texttt{O1} & \texttt{O2} & \texttt{O3} & \texttt{Os} & unroll & ffast & mcpu & ff+mcpu & +alias & lto \\
\midrule
\multicolumn{13}{l}{\textit{Instrukcje wektorowe (SIMD)}} \\[2pt]
\multirow{2}{*}{radix2\_v0}
  & GCC   & 0 & 0 & 0          & 0          & 0 & 0          & 0          & 0          & 0          & 0          & 0 \\
  & Clang & 0 & 0 & 0          & 0          & 0 & 0          & 0          & 0          & 0          & 0          & 0 \\
\addlinespace
\multirow{2}{*}{radix2\_v1}
  & GCC   & 0 & 0 & 5          & 5          & 0 & \textbf{15}& 0          & 0          & 0          & 0          & 5 \\
  & Clang & 0 & 0 & \textbf{9} & \textbf{9} & 0 & 9          & \textbf{8} & \textbf{9} & \textbf{8} & \textbf{8} & \textbf{9} \\
\addlinespace
\multirow{2}{*}{radix2\_v2\_soa}
  & GCC   & 0 & 0 & 0          & 0          & 0 & 0          & 0          & 0          & 0          & \textbf{16}& 0 \\
  & Clang & 0 & 0 & \textbf{9} & \textbf{9} & 0 & \textbf{9} & \textbf{8} & \textbf{9} & \textbf{8} & 8          & \textbf{9} \\
\addlinespace
\multirow{2}{*}{radix2\_v3\_neon}
  & GCC   & \textbf{9} & 8 & 8          & 8          & 8 & \textbf{32}& 8          & 8          & 8          & 16         & 8 \\
  & Clang & \textbf{9} & 8 & \textbf{17}& \textbf{17}& 8 & 17         & \textbf{16}& \textbf{17}& \textbf{16}& \textbf{16}& \textbf{17} \\
\addlinespace
\multirow{2}{*}{radix4}
  & GCC   & 0 & 0 & 0 & 0 & 0 & 0 & 0 & 0 & 0 & 0 & 0 \\
  & Clang & 0 & 0 & 0 & 0 & 0 & 0 & 0 & 0 & 0 & 0 & 0 \\
\addlinespace
\multirow{2}{*}{radix4\_naive}
  & GCC   & 0 & 0 & 0 & 0 & 0 & 0 & 0          & 0 & 0          & 0          & 0 \\
  & Clang & 0 & 0 & 0 & 0 & 0 & 0 & \textbf{31}& 0 & \textbf{31}& \textbf{31}& 0 \\
\addlinespace
\multirow{2}{*}{splitradix}
  & GCC   & 0 & 0 & \textbf{2} & \textbf{2} & 0 & \textbf{2} & \textbf{2} & \textbf{2} & \textbf{2} & \textbf{2} & \textbf{2} \\
  & Clang & 0 & 0 & 0          & 0          & 0 & 0          & 0          & 0          & 0          & 0          & 0 \\
\addlinespace
\multirow{2}{*}{stockham}
  & GCC   & 0 & 0 & 0          & 7          & 0 & \textbf{49}& 9          & 7          & 7          & 7          & 7 \\
  & Clang & 0 & 0 & \textbf{14}& \textbf{14}& \textbf{5} & 14         & \textbf{14}& \textbf{24}& \textbf{14}& \textbf{14}& \textbf{14} \\
\addlinespace
\multicolumn{13}{l}{\textit{Instrukcje skalarne}} \\[2pt]
\multirow{2}{*}{radix2\_v0}
  & GCC   & 11         & 17 & 13 & 13 & 13 & 36 & 13 & 13 & 13 & 13 & 13 \\
  & Clang & \textbf{9} & \textbf{13} & \textbf{13} & \textbf{13} & \textbf{13} & \textbf{13} & \textbf{13} & \textbf{13} & \textbf{13} & \textbf{13} & \textbf{13} \\
\addlinespace
\multirow{2}{*}{radix2\_v1}
  & GCC   & 10 & 10 & \textbf{1} & \textbf{1} & \textbf{8} & \textbf{3} & \textbf{8} & \textbf{8} & \textbf{8} & \textbf{8} & \textbf{1} \\
  & Clang & \textbf{8} & \textbf{8} & 8 & 8 & \textbf{8} & 8 & \textbf{8} & \textbf{8} & \textbf{8} & \textbf{8} & 8 \\
\addlinespace
\multirow{2}{*}{radix2\_v2\_soa}
  & GCC   & 10 & 10 & 8 & 8 & 8 & 24 & 8 & 8 & 8 & 16 & 8 \\
  & Clang & \textbf{8} & \textbf{8} & \textbf{8} & \textbf{8} & \textbf{8} & \textbf{8} & \textbf{8} & \textbf{8} & \textbf{8} & \textbf{8} & \textbf{8} \\
\addlinespace
\multirow{2}{*}{radix2\_v3\_neon}
  & GCC   & 10 & 10 & \textbf{8} & \textbf{24} & \textbf{8} & \textbf{24} & \textbf{24} & \textbf{24} & \textbf{24} & \textbf{32} & \textbf{24} \\
  & Clang & \textbf{8} & \textbf{8} & \textbf{8} & 32 & \textbf{8} & 32 & 32 & 32 & 32 & 32 & 32 \\
\addlinespace
\multirow{2}{*}{radix4}
  & GCC   & \textbf{22} & \textbf{34} & 28 & 28 & 28 & 28 & 28 & 28 & 28 & 28 & 28 \\
  & Clang & 20          & 28          & \textbf{28} & \textbf{28} & \textbf{28} & \textbf{28} & \textbf{28} & \textbf{28} & \textbf{28} & \textbf{28} & \textbf{28} \\
\addlinespace
\multirow{2}{*}{radix4\_naive}
  & GCC   & \textbf{33} & \textbf{40} & \textbf{34} & \textbf{34} & \textbf{34} & \textbf{34} & \textbf{34} & \textbf{34} & \textbf{34} & \textbf{34} & \textbf{34} \\
  & Clang & 31          & 34          & 34 & 34 & 34 & 34 & 34 & 34 & 34 & 34 & 34 \\
\addlinespace
\multirow{2}{*}{splitradix}
  & GCC   & \textbf{29} & \textbf{28} & \textbf{20} & \textbf{20} & \textbf{24} & \textbf{20} & \textbf{23} & \textbf{20} & \textbf{23} & \textbf{23} & \textbf{20} \\
  & Clang & 25          & 24          & 24          & 24          & 24          & 24          & 24          & 24          & 24          & 24          & 24 \\
\addlinespace
\multirow{2}{*}{stockham}
  & GCC   & \textbf{10} & \textbf{10} & 8 & 16 & 8 & 64 & 12 & 16 & 16 & 16 & 16 \\
  & Clang & 8           & 8           & \textbf{0} & \textbf{0} & \textbf{0} & \textbf{0} & \textbf{0} & \textbf{0} & \textbf{0} & \textbf{0} & \textbf{0} \\
\bottomrule
\end{tabular}
\end{table}

\subsection{Wpływ rozmiaru danych wejściowych}
Tutaj bedziemy mówić o chybieniach cache i o tym,że przy danych, które wychodza za cache duzo czasu zajmuje kopiowanie. 
Mozna jeszcze wsadzic tabelki 28 i 29, do zastanowienia. Jesli chce to sprawdzic to trzeba zobaczyc w jakim poziomie pamieci sa umieszczane dane,
bo czasy dostepu sa inne. Podobno mozna to sprawdzic z binarki dosc prosto, temat do przemyslenia.

%============================================================================
% Tabela 1.3   (etykieta: tab:med-gcc-1M)
% Mediana czasu wykonania transformaty dla GCC i 1048576 próbek, wersja sekwencyjna. Wartości w~milisekundach.
%============================================================================
\begin{table}[H]
\centering
\caption{Mediana czasu wykonania transformaty dla GCC i 1048576 próbek, wersja sekwencyjna. Wartości w~milisekundach.}
\label{tab:med-gcc-1M}
\scriptsize
\setlength{\tabcolsep}{3.0pt}
\begin{tabular}{lrrrrrrrrrrr}
\toprule
Algorytm & \texttt{O0} & \texttt{O1} & \texttt{O2} & \texttt{O3} & \texttt{Os} & unroll & ffast & mcpu & ff+mcpu & +alias & lto \\
\midrule
radix2\_v0 & 834,30 & 212,17 & 206,54 & 206,60 & 209,55 & 207,53 & 206,52 & 207,13 & 207,44 & 207,72 & 208,17 \\
radix2\_v1 & 638,30 & 211,61 & 208,53 & 208,09 & 212,02 & 207,28 & 209,20 & 207,36 & 209,96 & 210,51 & 212,75 \\
radix2\_v2\_soa & 473,97 & 211,86 & 209,70 & 208,33 & 207,83 & 212,03 & 216,72 & 208,50 & 216,78 & 215,87 & 212,16 \\
radix2\_v3\_neon & 419,65 & 212,35 & 210,82 & 210,87 & 214,28 & 208,78 & 212,63 & 211,66 & 213,72 & 214,83 & 212,40 \\
radix4 & 647,67 & 304,00 & 280,17 & 279,04 & 275,54 & 263,93 & 307,58 & 271,54 & 301,06 & 301,80 & 283,62 \\
radix4\_naive & 632,95 & 254,73 & 257,96 & 254,87 & 254,23 & 246,92 & 251,88 & 256,21 & 246,13 & 245,51 & 258,40 \\
splitradix & 745,81 & 634,16 & 597,59 & 596,89 & 602,57 & 601,41 & 630,62 & 553,78 & 605,62 & 600,52 & 600,61 \\
stockham & 310,10 & 171,62 & 170,54 & 91,42 & 167,89 & 92,67 & 90,69 & 91,51 & 91,50 & 91,78 & 91,27 \\
\addlinespace
fftw & 52,34 & 52,36 & 52,73 & 52,69 & 52,42 & 52,52 & 52,41 & 53,04 & 52,45 & 52,56 & 52,40 \\
\bottomrule
\end{tabular}
\end{table}



%============================================================================
% Tabela 1.33   (etykieta: tab:cache-gcc-1M)
% Współczynnik chybień pamięci podręcznej L1D i L2 dla GCC i 1048576 próbek, wersja sekwencyjna.
%============================================================================
\begin{table}[H]
\centering
\caption{Współczynnik chybień pamięci podręcznej L1D i L2 dla GCC i 1048576 próbek, wersja sekwencyjna.}
\label{tab:cache-gcc-1M}
\scriptsize
\setlength{\tabcolsep}{3.0pt}
\begin{tabular}{lrrrrrrrrrrr}
\toprule
Algorytm & \texttt{O0} & \texttt{O1} & \texttt{O2} & \texttt{O3} & \texttt{Os} & unroll & ffast & mcpu & ff+mcpu & +alias & lto \\
\midrule
\multicolumn{12}{l}{\textit{Chybienia L1D [\%]}} \\[2pt]
radix2\_v0 & 0,14 & 1,06 & 1,08 & 1,06 & 1,07 & 1,07 & 1,08 & 1,08 & 1,07 & 1,09 & 1,08 \\
radix2\_v1 & 0,26 & 1,12 & 1,94 & 2,01 & 1,10 & 2,00 & 1,11 & 1,14 & 1,12 & 1,14 & 2,17 \\
radix2\_v2\_soa & 0,52 & 2,09 & 2,15 & 1,98 & 2,13 & 1,87 & 1,90 & 1,98 & 1,84 & 4,52 & 2,53 \\
radix2\_v3\_neon & 0,85 & 4,76 & 5,43 & 4,87 & 4,94 & 5,03 & 5,03 & 5,08 & 5,21 & 5,09 & 6,03 \\
radix4 & 1,33 & 10,83 & 10,91 & 11,24 & 11,32 & 10,81 & 11,06 & 11,82 & 10,87 & 10,85 & 11,78 \\
radix4\_naive & 0,43 & 1,82 & 2,37 & 1,74 & 2,55 & 2,12 & 1,72 & 1,63 & 1,82 & 1,60 & 1,82 \\
splitradix & 2,69 & 18,06 & 19,10 & 18,97 & 18,93 & 18,97 & 18,70 & 18,86 & 18,88 & 19,04 & 19,07 \\
stockham & 0,76 & 3,24 & 3,18 & 1,51 & 5,63 & 1,57 & 1,47 & 1,44 & 1,57 & 1,39 & 1,41 \\
\addlinespace
\multicolumn{12}{l}{\textit{Chybienia L2D [\%]}} \\[2pt]
radix2\_v0 & 33,71 & 34,74 & 34,89 & 35,25 & 34,94 & 34,90 & 34,70 & 34,86 & 34,87 & 34,85 & 34,90 \\
radix2\_v1 & 31,46 & 31,38 & 31,15 & 30,77 & 31,53 & 31,03 & 31,34 & 30,99 & 31,49 & 31,29 & 30,76 \\
radix2\_v2\_soa & 26,39 & 30,00 & 29,04 & 29,29 & 29,25 & 30,86 & 30,95 & 30,38 & 31,90 & 30,56 & 26,82 \\
radix2\_v3\_neon & 24,48 & 31,44 & 29,53 & 31,23 & 31,08 & 30,34 & 31,02 & 30,72 & 30,62 & 30,54 & 28,10 \\
radix4 & 27,55 & 25,42 & 27,98 & 27,35 & 26,01 & 27,73 & 28,52 & 25,63 & 27,83 & 28,77 & 26,59 \\
radix4\_naive & 23,01 & 25,40 & 24,23 & 26,53 & 23,82 & 24,28 & 26,15 & 23,16 & 24,95 & 26,41 & 25,85 \\
splitradix & 27,73 & 28,77 & 26,75 & 25,19 & 25,56 & 26,01 & 26,94 & 26,17 & 25,72 & 25,48 & 28,34 \\
stockham & 28,89 & 32,79 & 33,09 & 26,19 & 32,65 & 26,24 & 26,20 & 26,29 & 26,03 & 26,24 & 26,34 \\
\bottomrule
\end{tabular}
\end{table}

%============================================================================
% Tabela 1.32   (etykieta: tab:cache-gcc-16k)
% Współczynnik chybień pamięci podręcznej L1D i L2 dla GCC i 16384 próbek, wersja sekwencyjna.
%============================================================================
\begin{table}[H]
\centering
\caption{Współczynnik chybień pamięci podręcznej L1D i L2 dla GCC i 16384 próbek, wersja sekwencyjna.}
\label{tab:cache-gcc-16k}
\scriptsize
\setlength{\tabcolsep}{3.0pt}
\begin{tabular}{lrrrrrrrrrrr}
\toprule
Algorytm & \texttt{O0} & \texttt{O1} & \texttt{O2} & \texttt{O3} & \texttt{Os} & unroll & ffast & mcpu & ff+mcpu & +alias & lto \\
\midrule
\multicolumn{12}{l}{\textit{Chybienia L1D [\%]}} \\[2pt]
radix2\_v0 & 0,19 & 1,44 & 1,42 & 1,46 & 1,42 & 1,48 & 1,43 & 1,47 & 1,45 & 1,38 & 1,47 \\
radix2\_v1 & 0,30 & 1,45 & 2,68 & 2,71 & 1,36 & 2,63 & 1,35 & 1,32 & 1,45 & 1,43 & 2,35 \\
radix2\_v2\_soa & 0,84 & 2,17 & 2,83 & 1,97 & 2,70 & 2,65 & 2,28 & 2,15 & 2,57 & 4,33 & 2,73 \\
radix2\_v3\_neon & 0,79 & 4,99 & 4,88 & 5,27 & 5,68 & 4,93 & 5,20 & 5,45 & 5,19 & 4,73 & 5,63 \\
radix4 & 0,83 & 6,59 & 7,51 & 6,95 & 6,87 & 6,80 & 6,96 & 7,29 & 7,25 & 6,29 & 6,94 \\
radix4\_naive & 0,24 & 1,33 & 1,44 & 1,96 & 1,48 & 1,20 & 1,49 & 1,05 & 1,54 & 1,57 & 1,64 \\
splitradix & 2,07 & 14,12 & 15,52 & 15,46 & 14,19 & 15,02 & 14,85 & 15,49 & 14,39 & 15,48 & 15,72 \\
stockham & 0,73 & 2,75 & 2,91 & 1,62 & 4,73 & 1,56 & 1,86 & 1,42 & 1,64 & 1,58 & 1,46 \\
\addlinespace
\multicolumn{12}{l}{\textit{Chybienia L2D [\%]}} \\[2pt]
radix2\_v0 & 0,07 & 0,23 & 0,22 & 0,00 & 0,16 & 0,14 & 0,04 & 0,07 & 0,24 & 0,26 & 0,22 \\
radix2\_v1 & 0,55 & 0,35 & 0,31 & 0,25 & 0,54 & 0,61 & 0,58 & 0,31 & 0,19 & 0,64 & 0,29 \\
radix2\_v2\_soa & 0,23 & 0,42 & 0,13 & 0,42 & 0,59 & 0,28 & 0,27 & 0,31 & 0,17 & 0,30 & 0,32 \\
radix2\_v3\_neon & 0,22 & 0,47 & 0,12 & 0,27 & 0,04 & 0,11 & 0,17 & 0,27 & 0,17 & 0,25 & 0,37 \\
radix4 & 0,18 & 0,16 & 0,17 & 0,27 & 0,18 & 0,12 & 0,11 & 0,20 & 0,08 & 0,05 & 0,47 \\
radix4\_naive & 0,17 & 0,10 & 0,10 & 0,11 & 0,30 & 0,15 & 0,12 & 0,17 & 0,12 & 0,16 & 0,28 \\
splitradix & 0,17 & 0,09 & 0,21 & 0,02 & 0,11 & 0,07 & 0,11 & 0,28 & 0,14 & 0,13 & 0,21 \\
stockham & 0,39 & 0,14 & 0,41 & 0,45 & 0,43 & 0,32 & 0,32 & 0,46 & 0,48 & 0,34 & 0,67 \\
\bottomrule
\end{tabular}
\end{table}



\section{Analiza wpływu optymalizacji programu wielowątkowego}

Ta sekcje trzeba przeprowadzic w bardzo analogiczny sposob jak te co jest wyzej, jak zrobisz sekcje sekwencyjna to to pojdzie w jeden wieczór


%============================================================================
% Tabela 1.68   (etykieta: tab:par-skal-gcc-64k)
% Przyspieszenie i efektywność zrównoleglenia względem jednego wątku dla GCC i 65536 próbek, wersja współbieżna.
%============================================================================
\begin{table}[H]
\centering
\caption{Przyspieszenie i efektywność zrównoleglenia względem jednego wątku dla GCC i 65536 próbek, wersja współbieżna.}
\label{tab:par-skal-gcc-64k}
\scriptsize
\setlength{\tabcolsep}{3.0pt}
\begin{tabular}{lrrrrrrrrrrr}
\toprule
Algorytm (wątki) & \texttt{O0} & \texttt{O1} & \texttt{O2} & \texttt{O3} & \texttt{Os} & unroll & ffast & mcpu & ff+mcpu & +alias & lto \\
\midrule
\multicolumn{12}{l}{\textit{Przyspieszenie względem jednego wątku [$\times$]}} \\[2pt]
radix2\_par (1) & 1,00 & 1,00 & 1,00 & 1,00 & 1,00 & 1,00 & 1,00 & 1,00 & 1,00 & 1,00 & 1,00 \\
radix2\_par (2) & 1,57 & 1,20 & 1,30 & 1,29 & 1,30 & 1,34 & 1,36 & 1,36 & 1,33 & 1,36 & 1,26 \\
radix2\_par (4) & 2,40 & 1,51 & 1,57 & 1,51 & 1,60 & 1,54 & 1,54 & 1,52 & 1,60 & 1,64 & 1,48 \\
radix2\_par (8) & 1,90 & 1,00 & 1,01 & 1,05 & 1,06 & 1,04 & 1,03 & 1,12 & 1,08 & 1,09 & 1,04 \\
\addlinespace
radix2\_par\_v2 (1) & 1,00 & 1,00 & 1,00 & 1,00 & 1,00 & 1,00 & 1,00 & 1,00 & 1,00 & 1,00 & 1,00 \\
radix2\_par\_v2 (2) & 1,83 & 1,44 & 1,38 & 1,41 & 1,47 & 1,36 & 1,46 & 1,49 & 1,49 & 1,46 & 1,38 \\
radix2\_par\_v2 (4) & 3,20 & 1,94 & 1,63 & 1,69 & 1,90 & 1,67 & 1,89 & 2,01 & 1,91 & 1,92 & 1,66 \\
radix2\_par\_v2 (8) & 2,54 & 1,25 & 1,10 & 1,12 & 1,27 & 1,11 & 1,24 & 1,34 & 1,28 & 1,29 & 1,11 \\
\addlinespace
radix4\_par (1) & 1,00 & 1,00 & 1,00 & 1,00 & 1,00 & 1,00 & 1,00 & 1,00 & 1,00 & 1,00 & 1,00 \\
radix4\_par (2) & 1,73 & 1,25 & 1,22 & 1,25 & 1,20 & 1,28 & 1,25 & 1,23 & 1,27 & 1,23 & 1,22 \\
radix4\_par (4) & 2,48 & 1,28 & 1,24 & 1,26 & 1,21 & 1,30 & 1,25 & 1,21 & 1,26 & 1,28 & 1,23 \\
radix4\_par (8) & 1,98 & 1,12 & 1,06 & 1,09 & 1,04 & 1,08 & 1,08 & 1,05 & 1,08 & 1,09 & 1,06 \\
\addlinespace
\multicolumn{12}{l}{\textit{Efektywność zrównoleglenia [\%]}} \\[2pt]
radix2\_par (1) & 100,0 & 100,0 & 100,0 & 100,0 & 100,0 & 100,0 & 100,0 & 100,0 & 100,0 & 100,0 & 100,0 \\
radix2\_par (2) & 78,3 & 60,2 & 65,2 & 64,6 & 64,9 & 66,8 & 68,0 & 68,1 & 66,3 & 67,9 & 62,8 \\
radix2\_par (4) & 60,1 & 37,7 & 39,4 & 37,9 & 39,9 & 38,5 & 38,6 & 38,1 & 39,9 & 40,9 & 37,0 \\
radix2\_par (8) & 23,8 & 12,5 & 12,6 & 13,2 & 13,3 & 13,0 & 12,9 & 14,0 & 13,5 & 13,7 & 13,0 \\
\addlinespace
radix2\_par\_v2 (1) & 100,0 & 100,0 & 100,0 & 100,0 & 100,0 & 100,0 & 100,0 & 100,0 & 100,0 & 100,0 & 100,0 \\
radix2\_par\_v2 (2) & 91,7 & 72,1 & 69,0 & 70,3 & 73,3 & 68,0 & 72,9 & 74,3 & 74,4 & 73,1 & 68,8 \\
radix2\_par\_v2 (4) & 80,0 & 48,5 & 40,9 & 42,3 & 47,4 & 41,7 & 47,2 & 50,2 & 47,7 & 48,0 & 41,5 \\
radix2\_par\_v2 (8) & 31,7 & 15,7 & 13,8 & 14,0 & 15,9 & 13,9 & 15,5 & 16,7 & 16,0 & 16,1 & 13,9 \\
\addlinespace
radix4\_par (1) & 100,0 & 100,0 & 100,0 & 100,0 & 100,0 & 100,0 & 100,0 & 100,0 & 100,0 & 100,0 & 100,0 \\
radix4\_par (2) & 86,6 & 62,4 & 61,2 & 62,5 & 60,0 & 63,8 & 62,7 & 61,7 & 63,7 & 61,7 & 60,8 \\
radix4\_par (4) & 61,9 & 32,0 & 30,9 & 31,5 & 30,3 & 32,6 & 31,3 & 30,3 & 31,6 & 31,9 & 30,8 \\
radix4\_par (8) & 24,7 & 14,0 & 13,3 & 13,7 & 13,0 & 13,5 & 13,5 & 13,2 & 13,6 & 13,6 & 13,2 \\
\bottomrule
\end{tabular}
\end{table}

%============================================================================
% Tabela 1.80   (etykieta: tab:par-cyk-instr-gcc-1M)
% Liczba cykli procesora, wykonanych instrukcji oraz IPC dla GCC i 1048576 próbek, wersja współbieżna. Cykle i~instrukcje podano w~milionach.
%============================================================================
\begin{table}[H]
\centering
\caption{Liczba cykli procesora, wykonanych instrukcji oraz IPC dla GCC i 1048576 próbek, wersja współbieżna. Cykle i~instrukcje podano w~milionach.}
\label{tab:par-cyk-instr-gcc-1M}
\scriptsize
\setlength{\tabcolsep}{3.0pt}
\begin{tabular}{lrrrrrrrrrrr}
\toprule
Algorytm (wątki) & \texttt{O0} & \texttt{O1} & \texttt{O2} & \texttt{O3} & \texttt{Os} & unroll & ffast & mcpu & ff+mcpu & +alias & lto \\
\midrule
\multicolumn{12}{l}{\textit{Cykle procesora}} \\[2pt]
radix2\_par (1) & 1149,1 & 392,9 & 388,3 & 384,7 & 398,5 & 399,8 & 393,6 & 386,4 & 393,4 & 403,1 & 389,7 \\
radix2\_par (2) & 1298,4 & 614,6 & 619,0 & 610,1 & 613,0 & 613,4 & 615,7 & 618,2 & 618,3 & 632,9 & 621,7 \\
radix2\_par (4) & 1352,3 & 1007,4 & 1017,9 & 980,8 & 1048,0 & 1017,3 & 1064,3 & 1008,3 & 901,0 & 1007,8 & 997,4 \\
radix2\_par (8) & 1425,6 & 1047,4 & 1035,6 & 1041,5 & 1075,4 & 1057,8 & 1084,5 & 1079,1 & 1051,1 & 1028,5 & 1052,4 \\
\addlinespace
radix2\_par\_v2 (1) & 1287,3 & 382,3 & 384,0 & 400,6 & 384,5 & 401,4 & 382,3 & 378,9 & 380,4 & 394,5 & 388,5 \\
radix2\_par\_v2 (2) & 1286,8 & 635,1 & 635,5 & 628,0 & 631,9 & 642,6 & 633,9 & 653,6 & 636,8 & 645,8 & 653,6 \\
radix2\_par\_v2 (4) & 1312,5 & 1153,7 & 1045,7 & 1059,4 & 1077,3 & 1082 & 1051,3 & 1056,1 & 1042,4 & 1083,7 & 1065,9 \\
radix2\_par\_v2 (8) & 1314,0 & 1003,7 & 1049,4 & 1018,4 & 1018,8 & 998,4 & 1020,5 & 1024,7 & 1016,3 & 1032,3 & 1029,4 \\
\addlinespace
radix4\_par (1) & 1100,7 & 542,5 & 512,8 & 515,0 & 536,9 & 492,7 & 551,6 & 509,8 & 560,7 & 559,5 & 515,5 \\
radix4\_par (2) & 1130,1 & 642,2 & 647,6 & 641,1 & 660,3 & 602,5 & 661,7 & 657,7 & 661,4 & 680,6 & 644,8 \\
radix4\_par (4) & 1216,0 & 924,1 & 916,4 & 1027,5 & 959,8 & 919,3 & 928,9 & 944,8 & 955,3 & 944,7 & 908,0 \\
radix4\_par (8) & 1258,2 & 915,1 & 915,0 & 1013,2 & 920,2 & 900,7 & 960,8 & 936,1 & 970,7 & 965,1 & 907,7 \\
\addlinespace
\multicolumn{12}{l}{\textit{Wykonane instrukcje}} \\[2pt]
radix2\_par (1) & 1345,3 & 332,3 & 224,3 & 224,2 & 315,5 & 226,3 & 308,2 & 308,6 & 308,7 & 308,7 & 224,2 \\
radix2\_par (2) & 1443,1 & 396,6 & 224,6 & 224,6 & 327,4 & 226,7 & 308,5 & 309,0 & 309,0 & 309,1 & 224,7 \\
radix2\_par (4) & 1443,6 & 397,3 & 225,2 & 225,2 & 328,1 & 227,4 & 309,2 & 309,8 & 309,6 & 309,7 & 225,3 \\
radix2\_par (8) & 1445,1 & 398,7 & 226,8 & 226,7 & 329,5 & 228,7 & 310,7 & 311,2 & 311,2 & 311,2 & 226,8 \\
\addlinespace
radix2\_par\_v2 (1) & 1425,1 & 398,3 & 260,0 & 260,0 & 388,8 & 249,4 & 375,2 & 375,7 & 375,8 & 375,8 & 261,0 \\
radix2\_par\_v2 (2) & 1425,4 & 398,8 & 260,4 & 260,3 & 389,3 & 250,0 & 375,7 & 376,3 & 376,3 & 376,2 & 261,5 \\
radix2\_par\_v2 (4) & 1425,7 & 399,4 & 261,0 & 261,0 & 390,1 & 250,8 & 376,3 & 376,8 & 377,0 & 377,1 & 262,1 \\
radix2\_par\_v2 (8) & 1427,3 & 400,9 & 262,3 & 262,3 & 391,3 & 251,9 & 377,6 & 378,1 & 378,2 & 378,2 & 263,4 \\
\addlinespace
radix4\_par (1) & 1080,4 & 307,3 & 284,1 & 284,0 & 291,4 & 257,9 & 294,6 & 297,7 & 308,2 & 308,2 & 284,1 \\
radix4\_par (2) & 1080,6 & 307,4 & 284,4 & 284,4 & 291,8 & 258,2 & 294,9 & 298,0 & 308,5 & 308,5 & 284,4 \\
radix4\_par (4) & 1080,9 & 307,9 & 284,9 & 284,9 & 292,2 & 258,6 & 295,3 & 298,5 & 309,0 & 309,2 & 284,8 \\
radix4\_par (8) & 1082,0 & 308,8 & 285,8 & 285,9 & 293,2 & 259,5 & 296,4 & 299,5 & 310,0 & 310,0 & 285,8 \\
\addlinespace
\multicolumn{12}{l}{\textit{Liczba instrukcji na cykl (IPC)}} \\[2pt]
radix2\_par (1) & 1,171 & 0,846 & 0,578 & 0,583 & 0,792 & 0,566 & 0,783 & 0,799 & 0,785 & 0,766 & 0,575 \\
radix2\_par (2) & 1,111 & 0,645 & 0,363 & 0,368 & 0,534 & 0,370 & 0,501 & 0,500 & 0,500 & 0,488 & 0,361 \\
radix2\_par (4) & 1,067 & 0,394 & 0,221 & 0,230 & 0,313 & 0,223 & 0,291 & 0,307 & 0,344 & 0,307 & 0,226 \\
radix2\_par (8) & 1,014 & 0,381 & 0,219 & 0,218 & 0,306 & 0,216 & 0,286 & 0,288 & 0,296 & 0,303 & 0,216 \\
\addlinespace
radix2\_par\_v2 (1) & 1,107 & 1,042 & 0,677 & 0,649 & 1,011 & 0,621 & 0,981 & 0,991 & 0,988 & 0,953 & 0,672 \\
radix2\_par\_v2 (2) & 1,108 & 0,628 & 0,410 & 0,415 & 0,616 & 0,389 & 0,593 & 0,576 & 0,591 & 0,583 & 0,400 \\
radix2\_par\_v2 (4) & 1,086 & 0,346 & 0,250 & 0,246 & 0,362 & 0,170 & 0,358 & 0,357 & 0,362 & 0,288 & 0,246 \\
radix2\_par\_v2 (8) & 1,086 & 0,399 & 0,250 & 0,258 & 0,384 & 0,252 & 0,370 & 0,369 & 0,372 & 0,366 & 0,256 \\
\addlinespace
radix4\_par (1) & 0,982 & 0,566 & 0,554 & 0,551 & 0,543 & 0,523 & 0,534 & 0,584 & 0,550 & 0,551 & 0,551 \\
radix4\_par (2) & 0,956 & 0,479 & 0,439 & 0,444 & 0,442 & 0,428 & 0,446 & 0,453 & 0,466 & 0,453 & 0,441 \\
radix4\_par (4) & 0,889 & 0,333 & 0,311 & 0,277 & 0,304 & 0,281 & 0,318 & 0,316 & 0,323 & 0,327 & 0,314 \\
radix4\_par (8) & 0,860 & 0,338 & 0,312 & 0,282 & 0,319 & 0,288 & 0,309 & 0,320 & 0,319 & 0,321 & 0,315 \\
\bottomrule
\end{tabular}
\end{table}

%============================================================================
% Tabela 1.57   (etykieta: tab:par-med-gcc-1M)
% Mediana czasu wykonania transformaty dla GCC i 1048576 próbek, wersja współbieżna. Wartości w~milisekundach.
%============================================================================
\begin{table}[H]
\centering
\caption{Mediana czasu wykonania transformaty dla GCC i 1048576 próbek, wersja współbieżna. Wartości w~milisekundach.}
\label{tab:par-med-gcc-1M}
\scriptsize
\setlength{\tabcolsep}{3.0pt}
\begin{tabular}{lrrrrrrrrrrr}
\toprule
Algorytm (wątki) & \texttt{O0} & \texttt{O1} & \texttt{O2} & \texttt{O3} & \texttt{Os} & unroll & ffast & mcpu & ff+mcpu & +alias & lto \\
\midrule
radix2\_par (1) & 637,42 & 211,43 & 209,69 & 211,15 & 212,29 & 212,82 & 209,59 & 211,09 & 209,59 & 211,45 & 210,77 \\
radix2\_par (2) & 415,02 & 197,09 & 196,34 & 196,99 & 198,30 & 196,15 & 199,48 & 197,21 & 195,69 & 196,75 & 197,49 \\
radix2\_par (4) & 288,64 & 182,77 & 182,70 & 181,09 & 183,14 & 182,01 & 181,97 & 182,29 & 182,42 & 182,42 & 181,37 \\
radix2\_par (8) & 312,27 & 204,87 & 203,18 & 202,12 & 204,69 & 201,27 & 202,64 & 204,96 & 204,12 & 203,77 & 202,07 \\
\addlinespace
radix2\_par\_v2 (1) & 709,30 & 213,34 & 210,20 & 210,46 & 211,23 & 209,57 & 208,53 & 210,26 & 208,95 & 210,87 & 210,78 \\
radix2\_par\_v2 (2) & 397,30 & 202,39 & 204,92 & 203,84 & 205,25 & 203,54 & 203,71 & 204,52 & 204,42 & 204,86 & 204,91 \\
radix2\_par\_v2 (4) & 245,13 & 179,38 & 178,99 & 178,37 & 181,66 & 179,40 & 179,22 & 179,37 & 180,69 & 181,30 & 179,91 \\
radix2\_par\_v2 (8) & 268,34 & 188,42 & 187,55 & 186,44 & 189,79 & 185,36 & 187,31 & 188,90 & 188,74 & 188,35 & 188,76 \\
\addlinespace
radix4\_par (1) & 654,86 & 301,52 & 279,47 & 283,32 & 294,55 & 270,43 & 298,29 & 279,77 & 307,36 & 306,00 & 283,99 \\
radix4\_par (2) & 383,40 & 217,00 & 211,01 & 215,11 & 218,61 & 195,09 & 217,98 & 216,45 & 221,98 & 223,15 & 210,99 \\
radix4\_par (4) & 276,50 & 182,69 & 182,53 & 183,74 & 187,90 & 166,68 & 184,60 & 186,89 & 188,05 & 189,99 & 180,58 \\
radix4\_par (8) & 286,94 & 190,93 & 189,87 & 192,42 & 193,49 & 177,01 & 194,49 & 190,89 & 200,64 & 198,21 & 187,17 \\
\bottomrule
\end{tabular}
\end{table}

%============================================================================
% Tabela 1.88   (etykieta: tab:par-ratio-gcc-1M)
% Stosunek czasu procesora do czasu ściany dla GCC i 1048576 próbek, wersja współbieżna. Wartości bezwymiarowe.
%============================================================================
\begin{table}[H]
\centering
\caption{Stosunek czasu procesora do czasu ściany dla GCC i 1048576 próbek, wersja współbieżna. Wartości bezwymiarowe.}
\label{tab:par-ratio-gcc-1M}
\scriptsize
\setlength{\tabcolsep}{3.0pt}
\begin{tabular}{lrrrrrrrrrrr}
\toprule
Algorytm (wątki) & \texttt{O0} & \texttt{O1} & \texttt{O2} & \texttt{O3} & \texttt{Os} & unroll & ffast & mcpu & ff+mcpu & +alias & lto \\
\midrule
radix2\_par (1) & 1,00 & 1,00 & 1,00 & 1,00 & 1,00 & 1,00 & 1,00 & 1,00 & 1,00 & 1,00 & 1,00 \\
radix2\_par (2) & 1,73 & 1,71 & 1,72 & 1,72 & 1,71 & 1,73 & 1,71 & 1,71 & 1,73 & 1,72 & 1,71 \\
radix2\_par (4) & 2,62 & 3,05 & 3,06 & 3,07 & 3,05 & 3,07 & 3,07 & 3,06 & 3,06 & 3,04 & 3,07 \\
radix2\_par (8) & 2,53 & 2,87 & 2,87 & 2,88 & 2,86 & 2,89 & 2,88 & 2,87 & 2,84 & 2,84 & 2,91 \\
\addlinespace
radix2\_par\_v2 (1) & 1,00 & 1,00 & 1,00 & 1,00 & 1,00 & 1,00 & 1,00 & 1,00 & 1,00 & 1,00 & 1,00 \\
radix2\_par\_v2 (2) & 1,79 & 1,71 & 1,71 & 1,72 & 1,72 & 1,73 & 1,72 & 1,72 & 1,71 & 1,71 & 1,71 \\
radix2\_par\_v2 (4) & 2,94 & 3,26 & 3,27 & 3,28 & 3,24 & 3,27 & 3,27 & 3,27 & 3,25 & 3,25 & 3,27 \\
radix2\_par\_v2 (8) & 2,71 & 3,00 & 3,02 & 3,01 & 3,00 & 3,05 & 3,01 & 3,04 & 3,02 & 3,04 & 2,99 \\
\addlinespace
radix4\_par (1) & 1,00 & 1,00 & 1,00 & 1,00 & 1,00 & 1,00 & 1,00 & 1,00 & 1,00 & 1,00 & 1,00 \\
radix4\_par (2) & 1,65 & 1,65 & 1,66 & 1,66 & 1,65 & 1,70 & 1,66 & 1,65 & 1,66 & 1,65 & 1,66 \\
radix4\_par (4) & 2,45 & 2,78 & 2,79 & 2,81 & 2,80 & 2,92 & 2,80 & 2,77 & 2,79 & 2,78 & 2,78 \\
radix4\_par (8) & 2,39 & 2,70 & 2,67 & 2,73 & 2,72 & 2,82 & 2,70 & 2,67 & 2,70 & 2,72 & 2,70 \\
\bottomrule
\end{tabular}
\end{table}

%============================================================================
% Tabela 1.69   (etykieta: tab:par-skal-gcc-1M)
% Przyspieszenie i efektywność zrównoleglenia względem jednego wątku dla GCC i 1048576 próbek, wersja współbieżna.
%============================================================================
\begin{table}[H]
\centering
\caption{Przyspieszenie i efektywność zrównoleglenia względem jednego wątku dla GCC i 1048576 próbek, wersja współbieżna.}
\label{tab:par-skal-gcc-1M}
\scriptsize
\setlength{\tabcolsep}{3.0pt}
\begin{tabular}{lrrrrrrrrrrr}
\toprule
Algorytm (wątki) & \texttt{O0} & \texttt{O1} & \texttt{O2} & \texttt{O3} & \texttt{Os} & unroll & ffast & mcpu & ff+mcpu & +alias & lto \\
\midrule
\multicolumn{12}{l}{\textit{Przyspieszenie względem jednego wątku [$\times$]}} \\[2pt]
radix2\_par (1) & 1,00 & 1,00 & 1,00 & 1,00 & 1,00 & 1,00 & 1,00 & 1,00 & 1,00 & 1,00 & 1,00 \\
radix2\_par (2) & 1,54 & 1,07 & 1,07 & 1,07 & 1,07 & 1,08 & 1,05 & 1,07 & 1,07 & 1,07 & 1,07 \\
radix2\_par (4) & 2,21 & 1,16 & 1,15 & 1,17 & 1,16 & 1,17 & 1,15 & 1,16 & 1,15 & 1,16 & 1,16 \\
radix2\_par (8) & 2,04 & 1,03 & 1,03 & 1,04 & 1,04 & 1,06 & 1,03 & 1,03 & 1,03 & 1,04 & 1,04 \\
\addlinespace
radix2\_par\_v2 (1) & 1,00 & 1,00 & 1,00 & 1,00 & 1,00 & 1,00 & 1,00 & 1,00 & 1,00 & 1,00 & 1,00 \\
radix2\_par\_v2 (2) & 1,79 & 1,05 & 1,03 & 1,03 & 1,03 & 1,03 & 1,02 & 1,03 & 1,02 & 1,03 & 1,03 \\
radix2\_par\_v2 (4) & 2,89 & 1,19 & 1,17 & 1,18 & 1,16 & 1,17 & 1,16 & 1,17 & 1,16 & 1,16 & 1,17 \\
radix2\_par\_v2 (8) & 2,64 & 1,13 & 1,12 & 1,13 & 1,11 & 1,13 & 1,11 & 1,11 & 1,11 & 1,12 & 1,12 \\
\addlinespace
radix4\_par (1) & 1,00 & 1,00 & 1,00 & 1,00 & 1,00 & 1,00 & 1,00 & 1,00 & 1,00 & 1,00 & 1,00 \\
radix4\_par (2) & 1,71 & 1,39 & 1,32 & 1,32 & 1,35 & 1,39 & 1,37 & 1,29 & 1,38 & 1,37 & 1,35 \\
radix4\_par (4) & 2,37 & 1,65 & 1,53 & 1,54 & 1,57 & 1,62 & 1,62 & 1,50 & 1,63 & 1,61 & 1,57 \\
radix4\_par (8) & 2,28 & 1,58 & 1,47 & 1,47 & 1,52 & 1,53 & 1,53 & 1,47 & 1,53 & 1,54 & 1,52 \\
\addlinespace
\multicolumn{12}{l}{\textit{Efektywność zrównoleglenia [\%]}} \\[2pt]
radix2\_par (1) & 100,0 & 100,0 & 100,0 & 100,0 & 100,0 & 100,0 & 100,0 & 100,0 & 100,0 & 100,0 & 100,0 \\
radix2\_par (2) & 76,8 & 53,6 & 53,4 & 53,6 & 53,5 & 54,2 & 52,5 & 53,5 & 53,6 & 53,7 & 53,4 \\
radix2\_par (4) & 55,2 & 28,9 & 28,7 & 29,2 & 29,0 & 29,2 & 28,8 & 28,9 & 28,7 & 29,0 & 29,1 \\
radix2\_par (8) & 25,5 & 12,9 & 12,9 & 13,1 & 13,0 & 13,2 & 12,9 & 12,9 & 12,8 & 13,0 & 13,0 \\
\addlinespace
radix2\_par\_v2 (1) & 100,0 & 100,0 & 100,0 & 100,0 & 100,0 & 100,0 & 100,0 & 100,0 & 100,0 & 100,0 & 100,0 \\
radix2\_par\_v2 (2) & 89,3 & 52,7 & 51,3 & 51,6 & 51,5 & 51,5 & 51,2 & 51,4 & 51,1 & 51,5 & 51,4 \\
radix2\_par\_v2 (4) & 72,3 & 29,7 & 29,4 & 29,5 & 29,1 & 29,2 & 29,1 & 29,3 & 28,9 & 29,1 & 29,3 \\
radix2\_par\_v2 (8) & 33,0 & 14,2 & 14,0 & 14,1 & 13,9 & 14,1 & 13,9 & 13,9 & 13,8 & 14,0 & 14,0 \\
\addlinespace
radix4\_par (1) & 100,0 & 100,0 & 100,0 & 100,0 & 100,0 & 100,0 & 100,0 & 100,0 & 100,0 & 100,0 & 100,0 \\
radix4\_par (2) & 85,4 & 69,5 & 66,2 & 65,9 & 67,4 & 69,3 & 68,4 & 64,6 & 69,2 & 68,6 & 67,3 \\
radix4\_par (4) & 59,2 & 41,3 & 38,3 & 38,6 & 39,2 & 40,6 & 40,4 & 37,4 & 40,9 & 40,3 & 39,3 \\
radix4\_par (8) & 28,5 & 19,7 & 18,4 & 18,4 & 19,0 & 19,1 & 19,2 & 18,3 & 19,1 & 19,3 & 19,0 \\
\bottomrule
\end{tabular}
\end{table}

%============================================================================
% Tabela 1.67   (etykieta: tab:par-skal-gcc-4k)
% Przyspieszenie i efektywność zrównoleglenia względem jednego wątku dla GCC i 4096 próbek, wersja współbieżna.
%============================================================================
\begin{table}[H]
\centering
\caption{Przyspieszenie i efektywność zrównoleglenia względem jednego wątku dla GCC i 4096 próbek, wersja współbieżna.}
\label{tab:par-skal-gcc-4k}
\scriptsize
\setlength{\tabcolsep}{3.0pt}
\begin{tabular}{lrrrrrrrrrrr}
\toprule
Algorytm (wątki) & \texttt{O0} & \texttt{O1} & \texttt{O2} & \texttt{O3} & \texttt{Os} & unroll & ffast & mcpu & ff+mcpu & +alias & lto \\
\midrule
\multicolumn{12}{l}{\textit{Przyspieszenie względem jednego wątku [$\times$]}} \\[2pt]
radix2\_par (1) & 1,00 & 1,00 & 1,00 & 1,00 & 1,00 & 1,00 & 1,00 & 1,00 & 1,00 & 1,00 & 1,00 \\
radix2\_par (2) & 1,30 & 0,59 & 0,58 & 0,58 & 0,63 & 0,59 & 0,61 & 0,63 & 0,61 & 0,61 & 0,57 \\
radix2\_par (4) & 1,63 & 0,47 & 0,45 & 0,44 & 0,51 & 0,47 & 0,48 & 0,48 & 0,49 & 0,49 & 0,44 \\
radix2\_par (8) & 0,94 & 0,22 & 0,20 & 0,19 & 0,23 & 0,20 & 0,21 & 0,22 & 0,21 & 0,21 & 0,19 \\
\addlinespace
radix2\_par\_v2 (1) & 1,00 & 1,00 & 1,00 & 1,00 & 1,00 & 1,00 & 1,00 & 1,00 & 1,00 & 1,00 & 1,00 \\
radix2\_par\_v2 (2) & 1,56 & 0,75 & 0,64 & 0,63 & 0,74 & 0,62 & 0,73 & 0,74 & 0,74 & 0,74 & 0,63 \\
radix2\_par\_v2 (4) & 2,17 & 0,61 & 0,51 & 0,50 & 0,62 & 0,49 & 0,58 & 0,62 & 0,60 & 0,60 & 0,50 \\
radix2\_par\_v2 (8) & 1,16 & 0,26 & 0,20 & 0,20 & 0,27 & 0,20 & 0,25 & 0,26 & 0,26 & 0,26 & 0,21 \\
\addlinespace
radix4\_par (1) & 1,00 & 1,00 & 1,00 & 1,00 & 1,00 & 1,00 & 1,00 & 1,00 & 1,00 & 1,00 & 1,00 \\
radix4\_par (2) & 1,36 & 0,73 & 0,71 & 0,73 & 0,73 & 0,71 & 0,74 & 0,73 & 0,73 & 0,75 & 0,70 \\
radix4\_par (4) & 1,61 & 0,54 & 0,54 & 0,57 & 0,55 & 0,54 & 0,57 & 0,58 & 0,60 & 0,57 & 0,54 \\
radix4\_par (8) & 0,88 & 0,25 & 0,24 & 0,25 & 0,26 & 0,24 & 0,26 & 0,26 & 0,27 & 0,26 & 0,25 \\
\addlinespace
\multicolumn{12}{l}{\textit{Efektywność zrównoleglenia [\%]}} \\[2pt]
radix2\_par (1) & 100,0 & 100,0 & 100,0 & 100,0 & 100,0 & 100,0 & 100,0 & 100,0 & 100,0 & 100,0 & 100,0 \\
radix2\_par (2) & 64,9 & 29,5 & 29,0 & 29,2 & 31,3 & 29,3 & 30,5 & 31,6 & 30,4 & 30,7 & 28,7 \\
radix2\_par (4) & 40,7 & 11,8 & 11,2 & 10,9 & 12,7 & 11,7 & 12,0 & 12,1 & 12,3 & 12,2 & 11,0 \\
radix2\_par (8) & 11,7 & 2,8 & 2,5 & 2,4 & 2,8 & 2,5 & 2,7 & 2,8 & 2,6 & 2,6 & 2,4 \\
\addlinespace
radix2\_par\_v2 (1) & 100,0 & 100,0 & 100,0 & 100,0 & 100,0 & 100,0 & 100,0 & 100,0 & 100,0 & 100,0 & 100,0 \\
radix2\_par\_v2 (2) & 77,8 & 37,5 & 31,8 & 31,3 & 37,2 & 31,0 & 36,7 & 37,1 & 37,0 & 37,1 & 31,7 \\
radix2\_par\_v2 (4) & 54,1 & 15,3 & 12,7 & 12,4 & 15,5 & 12,4 & 14,5 & 15,5 & 15,0 & 15,1 & 12,5 \\
radix2\_par\_v2 (8) & 14,5 & 3,3 & 2,6 & 2,5 & 3,3 & 2,5 & 3,1 & 3,3 & 3,3 & 3,2 & 2,6 \\
\addlinespace
radix4\_par (1) & 100,0 & 100,0 & 100,0 & 100,0 & 100,0 & 100,0 & 100,0 & 100,0 & 100,0 & 100,0 & 100,0 \\
radix4\_par (2) & 68,0 & 36,3 & 35,7 & 36,7 & 36,5 & 35,6 & 36,9 & 36,7 & 36,7 & 37,4 & 35,1 \\
radix4\_par (4) & 40,2 & 13,6 & 13,5 & 14,2 & 13,6 & 13,5 & 14,3 & 14,4 & 15,0 & 14,3 & 13,6 \\
radix4\_par (8) & 11,0 & 3,1 & 3,1 & 3,1 & 3,3 & 3,0 & 3,2 & 3,2 & 3,3 & 3,2 & 3,1 \\
\bottomrule
\end{tabular}
\end{table}


\subsection{Skalowanie z liczbą wątków}
\subsection{Wpływ flag kompilacji na warianty równoległe}